\begin{table}[!htbp]
\centering
\caption[Ciclul de viață runner plug-in]{Ciclul de viață al unui runner plug-in, de la upload până la
oprire. Stările sunt persistate în coloana \texttt{status} a tabelei
\texttt{runners}; tranzițiile sunt declanșate de API (creare/ștergere) și
de wrapper-ul din pod (marcarea \texttt{ready}).}
\label{tab:plugin_lifecycle}
\small
\begin{tabularx}{\columnwidth}{l X}
\toprule
\textbf{Stare} & \textbf{Trigger / descriere} \\
\midrule
\texttt{pending}  & INSERT inițial în DB; răspunsul HTTP returnează imediat. \\
\texttt{building} & API a creat Deployment-ul K8s prin \texttt{client-go}; pod-ul este \emph{Pending}. \\
\texttt{ready}    & Wrapper-ul din pod a citit manifestul din DB, a compilat scriptul și ascultă pe coadă. \\
\texttt{failed}   & Eroare la crearea Deployment-ului sau la compilarea scriptului; \texttt{status\_message} conține detalii. \\
\texttt{stopped}  & Oprire explicită; Deployment-ul este șters, dar rândul rămâne pentru audit (opțional). \\
\bottomrule
\end{tabularx}
\end{table}
